\documentclass[11pt,a4paper]{article}
\usepackage[margin=2.2cm]{geometry}
\usepackage{booktabs,multirow,longtable,graphicx,amsmath,hyperref,enumitem,xcolor}
\usepackage[T1]{fontenc}
\usepackage{lmodern}
\hypersetup{colorlinks=true,linkcolor=blue!60!black,urlcolor=blue!60!black}
\setlength{\parskip}{4pt}\setlength{\parindent}{0pt}
\newcommand{\tbl}[1]{\IfFileExists{tables/#1.tex}{\input{tables/#1.tex}}{\emph{(table #1 appears after the run)}}}

\title{Inverted Indexing, Sparse Retrieval, and the Vocabulary Mismatch Problem}
\author{ROLL1 \and ROLL2}
\date{5 September 2026}

\begin{document}
\maketitle

\section*{Datasets}
BEIR test splits loaded with \texttt{ir\_datasets} (\texttt{beir/<name>/test}). All runs on one laptop under WSL2
(Intel Core Ultra 7 155U, 7.6\,GB RAM for WSL, no GPU; Pyserini 2.3.0, Lucene 10, Java 21, PyTorch CPU).
\begin{itemize}[nosep]
  \item \textbf{SciFact}: 5\,183 documents, 300 test queries, 339 qrels. Standard split, no deviation.
  \item \textbf{FEVER}: 5\,416\,568 documents, 6\,666 test queries, 7\,937 qrels. Parts~1--3 on all queries; Parts~4--5 on a fixed random sample of 500 test queries (seed 0), because feedback and LLM/SPLADE encoding run per query on CPU.
  \item \textbf{HotpotQA}: 5\,233\,329 documents, 7\,405 test queries, 14\,810 qrels (two gold passages per question). Same 500-query sample for Parts~4--5.
\end{itemize}
Metrics: nDCG@10, Recall@100, MRR@10, MAP with \texttt{ir\_measures} on top-100 runs. Every table states the number of queries used, and each baseline is recomputed on the same queries as the method it is compared with.

\section{Part 1: Index}
One Lucene index per corpus with \texttt{python -m pyserini.index.lucene -collection JsonCollection -storeDocvectors}
(\texttt{contents} = title + text; default English analyzer; document vectors stored for Part~4). Build time is the wall-clock
time of the indexing process alone (corpus dump excluded), measured with no other job running.
\texttt{LuceneIndexReader} gives the statistics Part~4a needs: term-frequency vector (\texttt{get\_document\_vector}),
document length (its sum) and corpus document frequency (\texttt{get\_term\_counts}); Table~\ref{tab:p1demo} shows them for one gold document.

\tbl{part1}
\tbl{part1_demo}

\section{Part 2: Sparse retrieval baselines}
\texttt{LuceneSearcher} over each Part~1 index: (a) Pyserini default BM25 ($k_1{=}0.9$, $b{=}0.4$); (b) BM25 with $k_1, b$ chosen
per dataset by an exhaustive $5\times5$ grid over all test queries, selecting by nDCG@10 (Table~\ref{tab:grid}); (c) classic
TF-IDF by setting Lucene's \texttt{ClassicSimilarity} on the underlying \texttt{IndexSearcher}.

\tbl{part2}
\tbl{part2_grid}

PART2DISCUSSION

\section{Part 3: Vocabulary mismatch}
For every (query, gold document) pair, BM25 (tuned) \emph{fails} if the gold document is not in the top-10.
Lexical overlap = Jaccard coefficient between the analysed query tokens and the gold document's indexed term vector
(same analyzer as the index). Failures are bucketed as: \emph{no lexical overlap} (no shared analysed term);
\emph{abbreviation vs.\ expansion} (an all-caps token on one side whose initials appear as consecutive words on the other);
\emph{partial overlap / paraphrase} (fewer than half of the query terms in the gold document); \emph{high overlap but outranked}
(at least half of the query terms present, yet other documents score higher).

\tbl{part3_cats}
\tbl{part3_examples}
\tbl{part3_stats}
\tbl{part3_bins}
\tbl{part3_fig}

\paragraph{Verdict.} PART3VERDICT

\section{Part 4: Pseudo-relevance feedback}
\subsection{4a. Rocchio and RM3 (corpus-retrieved feedback)}
Own implementation (\texttt{rocchio.py}):
\[ w_t = \alpha\, f(q)[t] + \frac{\beta}{N}\sum_{d\in\text{top-}N} \tilde f(d)[t], \qquad \alpha=1,\ \beta=0.75, \]
with $f(q)$ and $\tilde f(d)$ length-normalised term frequencies ($\mathrm{tf}/|d|$). Feedback vectors are the Part~1 term vectors
of the top-$N$ tuned-BM25 documents (\texttt{IndexReader.get\_document\_vector}); terms in more than 10\% of the documents
(\texttt{get\_term\_counts}) are never added; the $k$ heaviest remaining centroid terms are kept. The re-weighted query is built
with Pyserini's query builder as a Boolean OR of \texttt{BoostQuery(TermQuery(contents:t), w\_t)} and re-searched.
RM3: \texttt{set\_rm3(fb\_terms=k, fb\_docs=N, original\_query\_weight=0.5)}.

\tbl{part4a}

\paragraph{Query drift.} Queries with the largest nDCG@10 drop after Rocchio, with the added terms:
\tbl{part4a_drift}
PART4ADRIFT

\subsection{4b. HyDE (LLM-generated feedback)}
Hypothetical documents are sampled locally on CPU with Qwen2.5-0.5B-Instruct (Transformers; $T{=}0.7$, top-$p$ 0.9;
two 96-token documents per query for SciFact, one 64-token document per query for the FEVER/HotpotQA samples), prompted to
write a passage that supports/refutes the claim or answers the question. The unchanged \texttt{rocchio.py} then receives the
analysed hypothetical documents as its feedback vectors. Variants on the same queries as 4a: naive concatenation
(query + hypothetical documents as one BM25 query), HyDE + Rocchio ($k\in\{10,20\}$), and 4a's corpus PRF.

\tbl{part4b}
\tbl{part4b_examples}

\paragraph{Feedback source vs.\ combination method.} PART4BVERDICT

\section{Part 5: SPLADE}
Checkpoint \texttt{naver/splade-cocondenser-ensembledistil} (SPLADE++ EnsembleDistil), $w_j=\max_i\log(1+\mathrm{ReLU}(o_{ij}))$.
SciFact is encoded on CPU with Pyserini's \texttt{SpladeDocumentEncoder} (max length 256, integer-quantised weights) and indexed
with \texttt{-impact -pretokenized}; the prebuilt SciFact index is run as a sanity check. Encoding 5M documents on CPU is not
feasible, so FEVER and HotpotQA use Pyserini's prebuilt \texttt{beir-v1.0.0-*.splade-pp-ed} impact indexes (same checkpoint);
queries are encoded locally with \texttt{SpladeQueryEncoder} and searched with \texttt{LuceneImpactSearcher}.

\tbl{part5}
\tbl{part5_index}

PART5DISCUSSION

\subsection*{Expansion terms: SPLADE vs.\ Rocchio vs.\ HyDE}
For 12 sample queries per dataset: the top-10 vocabulary terms SPLADE weights beyond the literal query word-pieces, next to
the top-10 Rocchio terms (4a, $N{=}10,k{=}20$) and the top-10 HyDE+Rocchio terms (4b, $k{=}20$). Overlap is counted after
stemming all three lists with the index analyzer.

\tbl{part5_overlap}
PART5OVERLAP
\tbl{part5_terms}

\section*{Use of LLM assistants}
Claude Code (Anthropic, Claude Fable 5.1) was used as a coding assistant inside VS Code/WSL: it wrote the scripts from our
instructions and the problem statement after inspecting the Pyserini API, ran them in our WSL environment, generated the LaTeX
tables from the result files and drafted this report, which we checked against the numbers and edited. The local LLM used
inside the experiments (Part~4b) is Qwen2.5-0.5B-Instruct, run offline. Shareable chat links: CHATLINKS.

\end{document}
